\section{Differential Cartier and the surviving residues}
\label{sec:cartier}

It is useful to normalize the same selected coefficient as a differential.
Over a perfect field $K$ of characteristic two, every rational function
has a unique expression $a(z)^2+zb(z)^2$: the elements $1,z$ form a
basis of $K(z)$ over $K(z^2)$, and perfectness supplies the square
roots of its coefficients. Define
$\mathcal C((a^2+zb^2)\,dz)=b\,dz$.
This is the differential Cartier operator. Multiplication by a square
pulls out its square root.

\begin{lemma}[Local rule]\label{lem:local}
For a fixed pole $\zeta$ and $b\in K$,
\[
 \mathcal C\!\left(\frac{b\,dz}{(z-\zeta)^{2j}}\right)=0,\qquad
 \mathcal C\!\left(\frac{b\,dz}{(z-\zeta)^{2j+1}}\right)
       =\frac{b^{1/2}\,dz}{(z-\zeta)^{j+1}}.
\]
The same definition gives
$\mathcal C(\sum a_nz^n\,dz)=\sum a_{2n+1}^{1/2}z^n\,dz$.
\end{lemma}
\begin{proof}
An even power of $h=z-\zeta$ is a square and contributes zero.
For odd powers, write $h^{2k+1}=(h^k)^2(z-\zeta)$.
The constant differential $\zeta\,dz$ has Cartier image zero,
whereas $\mathcal C(z\,dz)=dz$. Pulling out the square proves
the rule for all integer $k$, including negative $k$.
The coefficient rule follows by the same even/odd decomposition.
\end{proof}

For $e\ge1$ and $2^\ell\ge e$,
\begin{equation}\label{eq:selection}
 [z^0\,dz]\,\mathcal C^\ell\!
             \left(\frac{z^{e-1}}{A_d(z)}\,dz\right)
                   =p_d(2^\ell-e)\pmod2.
\end{equation}
Indeed $\ell$ odd selections choose original exponent $2^\ell-1$,
which in $z^{e-1}/A_d$ has coefficient $p_d(2^\ell-e)$.

\begin{lemma}[Removal of every higher pole]\label{lem:poles}
Put $\mu=\sum_{i=1}^d2^{\nu_2(i)}$ and
$a=\lceil\log_2\mu\rceil$. After $a$ Cartier applications to
a proper differential with denominator $A_d$, all surviving finite
poles are simple. Its simple residues subsequently undergo inverse
Frobenius at their fixed poles. No polynomial part is created.
\end{lemma}
\begin{proof}
The multiplicity at a root of odd order $q$ is
$\sum_{q\mid i,\ i\le d}2^{\nu_2(i)}\le\mu$.
By Lemma~\ref{lem:local}, a term of order $k>1$ survives $a$
applications only if $k-1$ is a positive multiple of $2^a$.
This is impossible for $k\le\mu\le2^a$. A simple residue is
replaced by its square root. The proper partial-fraction expansion
has no polynomial part, and every term in its Cartier image is again
proper, proving the final assertion.
\end{proof}

Here $\nu_2(i)$ denotes the exponent of two in the positive integer
$i$. For rational inputs over $\Ftwo$, the simple-pole tail is
eventually periodic: all its finitely many poles and residues lie
in one finite extension of $\Ftwo$, and inverse Frobenius is a
permutation of that finite field. This proves a finite cycle for
each supplied cutoff, with no cutoff-independent bound.

\subsection{Closure under both binary selectors}

The preceding lemma concerns repeated $\mathcal C$. To obtain a
complete observer for arbitrary remaining binary words, one must
also account for $\Lambda_0$.
For a rational function $R$ with coefficients in $\Ftwo$, define
\[
 \mathcal L_1(R\,dz)=\mathcal C(R\,dz),\qquad
 \mathcal L_0(R\,dz)=\mathcal C(zR\,dz).
 \tag{L}
\]
Their power-series expansions are respectively
$\Lambda_1(R)\,dz$ and $\Lambda_0(R)\,dz$, since square roots
fix the coefficients in $\Ftwo$.

\begin{lemma}[A uniform pole-order transient]\label{lem:binary-transient}
For a proper rational differential with denominator $A_d$,
each $\mathcal L_b$ preserves properness and introduces no
new finite pole. A pole of order $k$ acquires order at most
$\lceil k/2\rceil$. Consequently every binary word of length
$a=\lceil\log_2\mu\rceil$ reaches a proper differential
with only simple poles.
\end{lemma}
\begin{proof}
The assertion for $\mathcal L_1$ is the local Cartier rule.
At a pole $\zeta$, put $h=z-\zeta$. For $\mathcal L_0$,
write
\[
 \frac{bz}{h^k}\,dz
 =\frac{b}{h^{k-1}}\,dz+\frac{\zeta b}{h^k}\,dz.
\]
For $k=2j$ only the first summand survives Cartier, with
order $j$. For $k=2j+1$ only the second survives, with
order $j+1$. When $k=1$, the first summand is a constant
differential and is killed. Multiplication of a proper
rational function by $z$ adds at most a constant
polynomial part, which is likewise killed. Thus the
result is proper with the claimed poles and orders.

Iterating gives order at most $\lceil k/2^a\rceil$.
Every initial pole order is at most $\mu$, and $2^a\ge\mu$.
\end{proof}

Let $\mathcal R_d$ be the set of distinct roots of $A_d$ in
a finite splitting field $K$. These roots are nonzero and
the set is closed under Frobenius $\zeta\mapsto\zeta^2$.
Both the field and the root set can be computed by finite
polynomial factorization and adjoining roots. Define
\[
 \mathcal P_d=
 \left\{(b_\zeta)_{\zeta\in\mathcal R_d}\in K^{\mathcal R_d}:
              b_{\zeta^2}=b_\zeta^2\right\}.
 \tag{P}
\]

\begin{theorem}[A complete residue-profile observer]
\label{thm:residue-observer}
Proper rational differentials over $\Ftwo$ with simple poles
in $\mathcal R_d$ correspond bijectively to $\mathcal P_d$ by
\[
 R(z)\,dz=\sum_{\zeta\in\mathcal R_d}
                 \frac{b_\zeta}{z-\zeta}\,dz.
 \tag{PF}
\]
The profile space is finite and is closed under both binary
transitions. At each fixed pole their updates and the output are
\[
 \mathcal L_1:\ b_\zeta\longmapsto b_\zeta^{1/2},
 \qquad
 \mathcal L_0:\ b_\zeta\longmapsto(\zeta b_\zeta)^{1/2},
 \qquad
 \operatorname{out}(b)=\sum_{\zeta\in\mathcal R_d}
                             \frac{b_\zeta}{\zeta}\in\Ftwo.
 \tag{U}
\]
After the finite transient of
Lemma~\ref{lem:binary-transient}, this is an exact observer
for every remaining binary continuation.
\end{theorem}
\begin{proof}
A proper rational function with these simple poles has the
unique partial-fraction expansion (PF). Applying the field
Frobenius to its coefficients fixes the function precisely
when $b_{\zeta^2}=b_\zeta^2$. Conversely a compatible tuple
makes (PF) invariant under Frobenius, so its coefficients lie
in $\Ftwo$. There is no additional polynomial part.
This proves the bijection.

The local rule gives the first transition. For the second,
the decomposition
$z/(z-\zeta)=1+\zeta/(z-\zeta)$ leaves only the simple
fraction after Cartier, giving the second formula.
Each map preserves the compatibility equation, since
Frobenius and its inverse commute and
$(\zeta b_\zeta)^2=\zeta^2b_{\zeta^2}$.
They are $\Ftwo$-linear maps of the finite profile space.

The constant coefficient is $R(0)$, which is the displayed
sum because the characteristic is two. Frobenius permutes
its summands, so the sum lies in $\Ftwo$.
The profiles reconstruct the entire rational state, and
the two transition equations agree with the two binary
selectors by (L). Induction on the remaining word gives
the final assertion.
\end{proof}

If $S_d$ is the squarefree part of $A_d$, the profile space
is also the space of fractions $U/S_d$ with
$U\in\Ftwo[z]$, $\deg U<\deg S_d$. It therefore has
exactly $2^{\deg S_d}$ states. The initial transient can
be computed in $V_d$ by the numerator matrices and then
converted to this profile by partial fractions. Neither
representation suppresses a pole, a coefficient, or a
cancellation.

\begin{corollary}[Scalar output after repeated Cartier]
\label{cor:scalar-profile}
Let $\beta_\zeta$ denote the initial coefficient of
$(z-\zeta)^{-1}$ in the partial-fraction expansion of a
proper differential over $\Ftwo$ with denominator $A_d$.
For every $\ell\ge a$,
\[
 [z^0\,dz]\,\mathcal C^\ell(R\,dz)
   =\sum_{\zeta\in\mathcal R_d}
              \frac{\beta_\zeta^{\,2^{-\ell}}}{\zeta}.
 \tag{O}
\]
Here $2^{-\ell}$ denotes $\ell$ applications of inverse
Frobenius in $K$.
\end{corollary}
\begin{proof}
A higher pole never becomes simple under repeated
$\mathcal C$: while it survives its new order is
$1+(k-1)/2$, still greater than one. By time $a$ every
such term has vanished. The initial simple residues have
only undergone inverse Frobenius. Apply the output
formula (U).
\end{proof}

Grouping the terms in (O) by Frobenius orbit expresses the
answer as a sum of finite-field traces. If $\zeta$ has
orbit length $h$, compatibility puts its summand in
$\mathbb F_{2^h}$ and the orbit contributes
$\operatorname{Tr}_{\mathbb F_{2^h}/\Ftwo}
(\beta_\zeta^{2^{-\ell}}/\zeta)$.
This makes the scalar cancellation explicit.
A Frobenius-stable subfamily of poles has autonomous
updates under (U), but generally does not determine
the sum over the omitted poles.

\subsection{The actual observation lies beyond the transient}

For the linked values in \eqref{eq:delta}, $C=2^\ell$ with
$\ell\ge B+1$ by \eqref{eq:growth}. Since
$\mu\le d(1+\lfloor\log_2 d\rfloor)\le B(1+\log_2B)$,
we have $a\le B$ for $B\ge16$: use
$\log_2 B\le B/2$ and $B^2\le2^B$ in this range.
Also $B>2u$ gives $d=B-u+3>u$ and
$e-1=u-5<D$, so the actual initial numerator is proper.
Thus the actual selected observation
is beyond the higher-pole transient. This pays an actual parameter
bound; it is not an unspecified promise that enough iterations occur.
Together, \eqref{eq:selection} and (O) give the exact
scalar formula for that linked coefficient using all its
initial simple residues.

The complete state can nevertheless have a long simple-pole cycle.
Pole count alone cannot prove that. For example, $dS/S$ for a
squarefree polynomial $S$ has residue one at every simple pole and
is Cartier-fixed, however many poles $S$ has.

